% Options for packages loaded elsewhere
% Options for packages loaded elsewhere
\PassOptionsToPackage{unicode}{hyperref}
\PassOptionsToPackage{hyphens}{url}
\PassOptionsToPackage{dvipsnames,svgnames,x11names}{xcolor}
%
\documentclass[
  letterpaper,
  DIV=11,
  numbers=noendperiod]{scrartcl}
\usepackage{xcolor}
\usepackage{amsmath,amssymb}
\setcounter{secnumdepth}{-\maxdimen} % remove section numbering
\usepackage{iftex}
\ifPDFTeX
  \usepackage[T1]{fontenc}
  \usepackage[utf8]{inputenc}
  \usepackage{textcomp} % provide euro and other symbols
\else % if luatex or xetex
  \usepackage{unicode-math} % this also loads fontspec
  \defaultfontfeatures{Scale=MatchLowercase}
  \defaultfontfeatures[\rmfamily]{Ligatures=TeX,Scale=1}
\fi
\usepackage{lmodern}
\ifPDFTeX\else
  % xetex/luatex font selection
\fi
% Use upquote if available, for straight quotes in verbatim environments
\IfFileExists{upquote.sty}{\usepackage{upquote}}{}
\IfFileExists{microtype.sty}{% use microtype if available
  \usepackage[]{microtype}
  \UseMicrotypeSet[protrusion]{basicmath} % disable protrusion for tt fonts
}{}
\makeatletter
\@ifundefined{KOMAClassName}{% if non-KOMA class
  \IfFileExists{parskip.sty}{%
    \usepackage{parskip}
  }{% else
    \setlength{\parindent}{0pt}
    \setlength{\parskip}{6pt plus 2pt minus 1pt}}
}{% if KOMA class
  \KOMAoptions{parskip=half}}
\makeatother
% Make \paragraph and \subparagraph free-standing
\makeatletter
\ifx\paragraph\undefined\else
  \let\oldparagraph\paragraph
  \renewcommand{\paragraph}{
    \@ifstar
      \xxxParagraphStar
      \xxxParagraphNoStar
  }
  \newcommand{\xxxParagraphStar}[1]{\oldparagraph*{#1}\mbox{}}
  \newcommand{\xxxParagraphNoStar}[1]{\oldparagraph{#1}\mbox{}}
\fi
\ifx\subparagraph\undefined\else
  \let\oldsubparagraph\subparagraph
  \renewcommand{\subparagraph}{
    \@ifstar
      \xxxSubParagraphStar
      \xxxSubParagraphNoStar
  }
  \newcommand{\xxxSubParagraphStar}[1]{\oldsubparagraph*{#1}\mbox{}}
  \newcommand{\xxxSubParagraphNoStar}[1]{\oldsubparagraph{#1}\mbox{}}
\fi
\makeatother


\usepackage{longtable,booktabs,array}
\usepackage{calc} % for calculating minipage widths
% Correct order of tables after \paragraph or \subparagraph
\usepackage{etoolbox}
\makeatletter
\patchcmd\longtable{\par}{\if@noskipsec\mbox{}\fi\par}{}{}
\makeatother
% Allow footnotes in longtable head/foot
\IfFileExists{footnotehyper.sty}{\usepackage{footnotehyper}}{\usepackage{footnote}}
\makesavenoteenv{longtable}
\usepackage{graphicx}
\makeatletter
\newsavebox\pandoc@box
\newcommand*\pandocbounded[1]{% scales image to fit in text height/width
  \sbox\pandoc@box{#1}%
  \Gscale@div\@tempa{\textheight}{\dimexpr\ht\pandoc@box+\dp\pandoc@box\relax}%
  \Gscale@div\@tempb{\linewidth}{\wd\pandoc@box}%
  \ifdim\@tempb\p@<\@tempa\p@\let\@tempa\@tempb\fi% select the smaller of both
  \ifdim\@tempa\p@<\p@\scalebox{\@tempa}{\usebox\pandoc@box}%
  \else\usebox{\pandoc@box}%
  \fi%
}
% Set default figure placement to htbp
\def\fps@figure{htbp}
\makeatother


% definitions for citeproc citations
\NewDocumentCommand\citeproctext{}{}
\NewDocumentCommand\citeproc{mm}{%
  \begingroup\def\citeproctext{#2}\cite{#1}\endgroup}
\makeatletter
 % allow citations to break across lines
 \let\@cite@ofmt\@firstofone
 % avoid brackets around text for \cite:
 \def\@biblabel#1{}
 \def\@cite#1#2{{#1\if@tempswa , #2\fi}}
\makeatother
\newlength{\cslhangindent}
\setlength{\cslhangindent}{1.5em}
\newlength{\csllabelwidth}
\setlength{\csllabelwidth}{3em}
\newenvironment{CSLReferences}[2] % #1 hanging-indent, #2 entry-spacing
 {\begin{list}{}{%
  \setlength{\itemindent}{0pt}
  \setlength{\leftmargin}{0pt}
  \setlength{\parsep}{0pt}
  % turn on hanging indent if param 1 is 1
  \ifodd #1
   \setlength{\leftmargin}{\cslhangindent}
   \setlength{\itemindent}{-1\cslhangindent}
  \fi
  % set entry spacing
  \setlength{\itemsep}{#2\baselineskip}}}
 {\end{list}}
\usepackage{calc}
\newcommand{\CSLBlock}[1]{\hfill\break\parbox[t]{\linewidth}{\strut\ignorespaces#1\strut}}
\newcommand{\CSLLeftMargin}[1]{\parbox[t]{\csllabelwidth}{\strut#1\strut}}
\newcommand{\CSLRightInline}[1]{\parbox[t]{\linewidth - \csllabelwidth}{\strut#1\strut}}
\newcommand{\CSLIndent}[1]{\hspace{\cslhangindent}#1}



\setlength{\emergencystretch}{3em} % prevent overfull lines

\providecommand{\tightlist}{%
  \setlength{\itemsep}{0pt}\setlength{\parskip}{0pt}}



 


% ============================================================
% loani.tex  —  include-in-header snippet, LuaLaTeX
% IPA handled by ipa.lua Lua filter at AST level.
% No ucharclasses, no babel, no runtime font switching.
% ============================================================

\usepackage{iftex}
\RequireLuaTeX

\usepackage{fontspec}

\setmainfont{Georgia}[
  Ligatures = TeX
]
% \usepackage[hebrew]{babel}
% Hebrew font — switched via ucharclasses (XeTeX) or
% directly referenced where needed
% \newfontfamily\hebrewfont{Alef}[
%   Script = Hebrew
% ]
% \newfontfamily\babelfont{Alef}[
%   Script = Hebrew
% ]
% \usepackage[hebrew, english]{babel}
% \babelfont[hebrew]{rm}{Alef}
% \usepackage[english, bidi=basic]{babel}
% \usepackage[english, hebrew, bidi=basic]{babel}
\usepackage[english, bidi=basic]{babel}

% \babelprovide[onchar=fonts ids, main=false]{hebrew} % all RTL, IPA wrongs
% \babelprovide[onchar=ids fonts, main=false]{english}
% \babelprovide[onchar=fonts ids, main=true]{english}
\babelprovide[onchar=fonts ids]{hebrew}

\babelfont[hebrew]{rm}{Alef}
% IPA font — called explicitly by ipa.lua filter output
\newfontfamily\ipafont{Source Sans 3}

% ------------------------------------------------------------
% Hebrew bidi — luatexja-style or direct LuaTeX primitive.
% Since babel is not loaded, we use the LuaTeX direction
% primitive directly: TRT = right-to-left, TLT = left-to-right
% Hebrew Unicode chars get RTL from the UBA via luaotfload.
% ------------------------------------------------------------
\directlua{
  luatexbase.add_to_callback("pre_linebreak_filter",
    function(head) return head end,
    "dummy_bidi")
}

\usepackage[hyperfootnotes=false]{hyperref}
% ============================================================
% header.tex  —  include-in-header snippet
% Replaces $var$ Pandoc syntax with \ifdefined\doc... guards.
% Vars injected via vars.lua Lua filter from YAML front matter:
%   \doctitle      \docsubtitle   \docauthor   \docdate
%   \docmeta       \docfoot       \doclogoleft \doclogoright
%
% No fancyhdr. Header fires via \AtBeginDocument,
% footer via \AtEndDocument — bidi-safe.
% ============================================================

\usepackage{graphicx}
\usepackage{tabularx}
\usepackage{xcolor}
\usepackage{array}
\usepackage{tikz}

% ---------- header font (adjust to taste) ----------
\newfontfamily\headerfont{Source Sans 3}

% ---------- footer colors ----------
\definecolor{footerbg}{HTML}{000000}
\definecolor{footerfg}{HTML}{FFFFFF}

% ---------- column types ----------
\newcolumntype{L}{>{\raggedright\arraybackslash}m{10mm}}
\newcolumntype{C}{>{\centering\arraybackslash}X}
\newcolumntype{R}{>{\raggedleft\arraybackslash}m{45mm}}

% vertical separator rule
\newcommand{\vsep}{\color{black}\vrule width 1pt}

% ============================================================
% CUSTOM HEADER
% ============================================================
\newcommand{\CustomHeader}{%
% \LTR{%
\noindent
\begin{tabularx}{\textwidth}{L !{\vsep} C !{\vsep} R}

% LEFT LOGO
\ifdefined\doclogoleft
  \includegraphics[width=10mm,height=10mm,keepaspectratio]{\doclogoleft}%
\fi
&
% CENTER: title + subtitle stacked in a minipage so \\ is a
% line break within the cell, not a new table row
\begin{minipage}[c]{\linewidth}
\centering
\ifdefined\doctitle
  {\Large\bfseries\headerfont\doctitle}%
\fi
\ifdefined\docsubtitle
  \\[2pt]{\normalsize\headerfont\docsubtitle}%
\fi
\end{minipage}
% &
% % CENTER TITLE BLOCK
% \ifdefined\doctitle
%   {\Large\bfseries\doctitle}%
% \fi
% \ifdefined\docsubtitle
%   \\\normalsize\docsubtitle
% \fi
&
% RIGHT LOGO
\ifdefined\doclogoright
  \includegraphics[height=10mm,keepaspectratio]{\doclogoright}%
\fi

\end{tabularx}

\hrule
% \vspace{6pt}

% ---------- META AREA ----------
{\small
\ifdefined\docauthor
  \textbf{Author:} \docauthor\\
\fi
\ifdefined\docdate
  \textbf{Date:} \docdate
\fi
\ifdefined\docmeta
  \\\docmeta\\
  snc: 16132.  
  {\ipafont ɪ ʃ ɔ ː ʁ}  
  {\ipafont wer es [pɛʟɛχɪʟ] ausspricht, und nicht [pɛʀɛχɪʟ]}
\fi
}
\vspace{6pt}

\hrule

% \vspace{12pt}
% }% end \LTR
}% end \CustomHeader

% ============================================================
% CUSTOM FOOTER  (overlay, anchored to page bottom)
% ============================================================
\newcommand{\CustomFooter}{%
\begin{tikzpicture}[remember picture, overlay]
\node[
  anchor=south,
  draw,
  fill=footerbg,
  text=footerfg,
  inner sep=10pt,
  rounded corners=2pt,
  text width=\dimexpr\textwidth-20pt\relax
]
at ([yshift=5mm]current page.south)
{%
  \begin{center}
  {\small
  \ifdefined\docfoot
    % {\docfoot}%
    {\headerfont\docfoot}%
  \fi
  }
  \end{center}
};
\end{tikzpicture}%
}% end \CustomFooter

% ============================================================
% FIRE header at doc start, footer at doc end
% ============================================================
% \AtBeginDocument{\CustomHeader}
\AtEndDocument{\CustomFooter}
\renewcommand{\maketitle}{}
\KOMAoption{captions}{tableheading}
\usepackage{fancyvrb}
\usepackage{fvextra}
\RecustomVerbatimEnvironment{verbatim}{Verbatim}{breaklines=true,breakanywhere=true,fontsize=\small}
\newcommand{\doctitle}{gemini paper}
\newcommand{\docdate}{2026-04-01}
\newcommand{\docmeta}{class: [16838] germanische sprachen im vergleich tutor: hüning/konvicka term: WS25/26}
\newcommand{\doclogoleft}{../assets/fadenlogo-fwb.png}
\newcommand{\doclogoright}{../assets/dcl-head-niemand.png}
\makeatletter
\@ifpackageloaded{caption}{}{\usepackage{caption}}
\AtBeginDocument{%
\ifdefined\contentsname
  \renewcommand*\contentsname{Table of contents}
\else
  \newcommand\contentsname{Table of contents}
\fi
\ifdefined\listfigurename
  \renewcommand*\listfigurename{List of Figures}
\else
  \newcommand\listfigurename{List of Figures}
\fi
\ifdefined\listtablename
  \renewcommand*\listtablename{List of Tables}
\else
  \newcommand\listtablename{List of Tables}
\fi
\ifdefined\figurename
  \renewcommand*\figurename{Figure}
\else
  \newcommand\figurename{Figure}
\fi
\ifdefined\tablename
  \renewcommand*\tablename{Table}
\else
  \newcommand\tablename{Table}
\fi
}
\@ifpackageloaded{float}{}{\usepackage{float}}
\floatstyle{ruled}
\@ifundefined{c@chapter}{\newfloat{codelisting}{h}{lop}}{\newfloat{codelisting}{h}{lop}[chapter]}
\floatname{codelisting}{Listing}
\newcommand*\listoflistings{\listof{codelisting}{List of Listings}}
\makeatother
\makeatletter
\makeatother
\makeatletter
\@ifpackageloaded{caption}{}{\usepackage{caption}}
\@ifpackageloaded{subcaption}{}{\usepackage{subcaption}}
\makeatother
\usepackage{bookmark}
\IfFileExists{xurl.sty}{\usepackage{xurl}}{} % add URL line breaks if available
\urlstyle{same}
\hypersetup{
  colorlinks=true,
  linkcolor={blue},
  filecolor={Maroon},
  citecolor={Blue},
  urlcolor={Blue},
  pdfcreator={LaTeX via pandoc}}


\title{gemini paper}
\author{}
\date{2026-04-01}
\begin{document}
\maketitle

\CustomHeader


\section{index}\label{index}

\subsection{snc}\label{snc}

\begin{itemize}
\tightlist
\item
  16144.1.pdf
\item
  wks. no fro q.yml
\item
  no.
\end{itemize}

\subsubsection{16145.}\label{section}

\begin{itemize}
\tightlist
\item
  from parent qyml
\end{itemize}

\section{the einleitung}\label{the-einleitung}

inspired by the paper \emph{Empirical evidence of Large Language Model's
influence on human spoken communication}, Yakura et al. (2025), who
indeed found (evidence) for GPT influenced human language after the
introduction of chatGPT we tried to replicate the pipeline of building
an AI vocabulary (gpt preferred lemmata) and compare frequencies of
gpt-typical words across pre- and post chatGPT human language corpora.
The first draft essai proves their hypothesis that LLM generated
language manifests within human natural language.

the embedding of that investigation into the context of the class
subject \emph{germanische sprachen im vergleich} is still due; first
idea is the projection of the Yakura et al. (2025) findings onto a
german language corpus and see if these are still valid although that
may rather be a pragmatics investigation.

\subsection{preliminary}\label{preliminary}

Our findings are still limited to a yet very small corpus of texts after
the introduction of the google gemini chat agent to the german public in
03/2024, cf. Wikipedia and Google (2026). In contrast to Yakura et al.
(2025) and out of resources considerations we decided for gemini as
basis for our AI generated vocabulary and for another text corpus
(german bundestag plenary protocols, DIP (2026)) than youtube/podcast
audio for the same reasons. That limits our post-AI corpus to a small
timeframe between 03/2024 up to now. With expanding that corpus to a
wider spectrum with including other sources we may harden our results.

\section{hypothesis}\label{hypothesis}

following Yakura et al. (2025) we assumed that the consuming of LLM
generated language influences the human production of language such that
vocabulary typical for LLM output will be found with higher frequencies
in human language corpora dating after chat agent introduction.

\section{methods}\label{methods}

\subsection{snc}\label{snc-1}

16062.1.2.16063.1

please cf. Schwarz (2026) for the corpus building and evaluation script
(still messy.)

\subsection{data}\label{data}

our human langugae data consists of raw texts from german bundestag
plenary protocols (DIP (2026)). the LLM corpus consists of model
summaries of a first subset of these texts generated with the following
prompt: Section~\ref{sec-gemini-prompt}.

\subsubsection{corpus subsets}\label{corpus-subsets}

\begin{longtable}[]{@{}lr@{}}
\toprule\noalign{}
target & tokens \\
\midrule\noalign{}
\endhead
\bottomrule\noalign{}
\endlastfoot
gemini & 3866 \\
human-pre & 1640446 \\
human-post & 1940385 \\
\end{longtable}

\subsubsection{gemini prompt}\label{sec-gemini-prompt}

\begin{verbatim}
[1] "System prompt: "                                                                                                                                                                                                                                                                                                                                                                                                                                                                                                                                                                           
[2] "You are a member of german parliament. Prepare a summary of the text provided to present at a local community meeting of your party members. Output in german language, no preamble, no extra information, just the plain text. Wordcount maximal 300 words, containing not more than 5% of the keywords of the text provided and explicitly not just a list of keywords but an entertaining text. You are supposed to interprete freely, including background insights on daily politics. Keep in mind thatthe text will be used as is as keynotes to the talk being held to the locals. "
[3] "Text:"                                                                                                                                                                                                                                                                                                                                                                                                                                                                                                                                                                                     
\end{verbatim}

\subsection{computation}\label{computation}

we first devised AI-typical lemmata in the model corpus which are
distinct for that corpus by log frequency calculation, see
Figure~\ref{fig-gpt.dplot}

\begin{figure}[H]

\centering{

\pandocbounded{\includegraphics[keepaspectratio]{gemini_files/figure-pdf/fig-gpt.dplot-1.pdf}}

}

\caption{\label{fig-gpt.dplot}lemma gpt scores over targets}

\end{figure}%

\section{evaluation}\label{evaluation}

\subsection{basic descriptive}\label{basic-descriptive}

to first gather an insight, yet with simple descriptive stats comparing
the raw frequencies of gpt-preferred lemmas in pre- and post-gemini
onset we find that in the target corpus the occurences of these lemma
increase, only by small amount (see Table~\ref{tbl-s.df}) and hard to
visualise (see Figure~\ref{fig-boxplot01}). wether these findings become
relevant, we'll see in Section~\ref{sec-lm-1} where we evaluate the
frequencies with a linear regression model.

\begin{longtable}[]{@{}lr@{}}

\caption{\label{tbl-s.df}GPT lemma frequencies (table) over target.
(freq / vH)}

\tabularnewline

\toprule\noalign{}
target & freq \\
\midrule\noalign{}
\endhead
\bottomrule\noalign{}
\endlastfoot
pre & 0.4384 \\
post & 0.4414 \\
DIFF: & 0.0030 \\

\end{longtable}

\begin{figure}[H]

\centering{

\pandocbounded{\includegraphics[keepaspectratio]{gemini_files/figure-pdf/fig-boxplot01-1.pdf}}

}

\caption{\label{fig-boxplot01}GPT lemma frequencies (boxplot) over
target. (freq / vH)}

\end{figure}%

\subsection{responsible lemmata}\label{responsible-lemmata}

The 9 lemma where 1. the relative frequency of lemma in target=post is
higher than in target=pre (n=27) and 2. that frequency exceeds the mean
relative frequency is displayed in Figure~\ref{fig-barplot}.

\begin{figure}[H]

\centering{

\pandocbounded{\includegraphics[keepaspectratio]{gemini_files/figure-pdf/fig-barplot-1.pdf}}

}

\caption{\label{fig-barplot}lemma frequency increase of relevant AI
typical lemma}

\end{figure}%

\subsubsection{lemma descriptive output}\label{lemma-descriptive-output}

\subsubsection{lemma linear model
output}\label{lemma-linear-model-output}

\subsection{linear regression}\label{sec-lm-1}

to prove descriptive results, we compute the stability of the frequency
increase for target- vs.~reference corpus with a linear regression model
using R's lme4::lmer() function, cf. Bates et al. (2015). coefficents
are printed below, where \texttt{freq\_pmw} is the lemma frequencies/pmw
over corpus; \texttt{post} defines reference resp. target corpus
{[}post==TRUE\textbar FALSE{]} (pre/post) and gp as numerical variable
representing the gpt-score of the corresponding lemma i.e.~wether it
scores high (positive values) or low (negative values) in terms of being
preferredly used by the chat agent.

\subsubsection{basic (lm)}\label{basic-lm}

formula: \texttt{frequency.pmw\ \textasciitilde{}\ post\ +\ gpt.score}

\begin{verbatim}

Call:
lm(formula = freq_pmw ~ post + gp, data = df_lg.norm)

Residuals:
   Min     1Q Median     3Q    Max 
 -1400   -435   -388   -243 436922 

Coefficients:
            Estimate Std. Error t value Pr(>|t|)    
(Intercept)   833.75      38.64  21.576  < 2e-16 ***
postTRUE       83.94      20.64   4.067 4.77e-05 ***
gp             68.11       5.20  13.097  < 2e-16 ***
---
Signif. codes:  0 '***' 0.001 '**' 0.01 '*' 0.05 '.' 0.1 ' ' 1

Residual standard error: 3765 on 133087 degrees of freedom
  (1393 observations deleted due to missingness)
Multiple R-squared:  0.0014,    Adjusted R-squared:  0.001385 
F-statistic: 93.32 on 2 and 133087 DF,  p-value: < 2.2e-16
\end{verbatim}

\subsubsection{random effects model
(lmer)}\label{random-effects-model-lmer}

formula:
\texttt{frequency.pmw\ \textasciitilde{}\ post\ +\ gpt.score\ +\ (1\textbar{}lemma)}

\begin{verbatim}
Linear mixed model fit by REML. t-tests use Satterthwaite's method [
lmerModLmerTest]
Formula: freq_pmw ~ post + gp + (1 | lemma)
   Data: df_lg.norm

REML criterion at convergence: 2490672

Scaled residuals: 
    Min      1Q  Median      3Q     Max 
-48.624  -0.047  -0.032   0.008  74.360 

Random effects:
 Groups   Name        Variance Std.Dev.
 lemma    (Intercept) 8839356  2973    
 Residual             1754704  1325    
Number of obs: 133090, groups:  lemma, 99649

Fixed effects:
             Estimate Std. Error        df t value Pr(>|t|)    
(Intercept) 4.344e+02  3.966e+01 1.055e+05  10.953   <2e-16 ***
postTRUE    1.246e+02  9.496e+00 5.602e+04  13.126   <2e-16 ***
gp          2.932e+01  5.723e+00 1.017e+05   5.124    3e-07 ***
---
Signif. codes:  0 '***' 0.001 '**' 0.01 '*' 0.05 '.' 0.1 ' ' 1

Correlation of Fixed Effects:
         (Intr) psTRUE
postTRUE -0.110       
gp        0.959  0.009
\end{verbatim}

\begin{verbatim}
3 x 3 Matrix of class "corMatrix"
            (Intercept)     postTRUE          gp
(Intercept)   1.0000000 -0.110159806 0.959345449
postTRUE     -0.1101598  1.000000000 0.008711618
gp            0.9593454  0.008711618 1.000000000
\end{verbatim}

\paragraph{helper interpretation, to be
tested}\label{helper-interpretation-to-be-tested}

the coefficients interesting for us are the gp and post estimates. here
we test the association between the gpt score of a lemma and its
estimated frequency and its showing that a general increase of frequency
is estimated if the score rises (=the lemma is within the lemmas
preferred used by gemini) and that for the post-gpt corpus this increase
(\texttt{6.75072\%}) is significant for single lemmata (and not random
in data).

in the fixed effects correlation output of the lmer() model we see that
the gpt score correlates with the target corpus frequency for lemma by
\texttt{1}.

\subsubsection{anova of mixed effects model
{[}out{]}}\label{anova-of-mixed-effects-model-out}

\section*{references}\label{references}
\addcontentsline{toc}{section}{references}

\phantomsection\label{refs}
\begin{CSLReferences}{1}{0}
\bibitem[\citeproctext]{ref-bates_fitting_2015}
Bates, Douglas, Martin Mächler, Ben Bolker, and Steve Walker. 2015.
{``Fitting {Linear} {Mixed}-{Effects} {Models} {Using} Lme4.''}
\emph{Journal of Statistical Software} 67 (1): 1--48.
\url{https://doi.org/10.18637/jss.v067.i01}.

\bibitem[\citeproctext]{ref-dip_dip_2026}
DIP. 2026. {``{DIP} - {Bundestagsprotokolle}.''} Docs. \emph{DIP - API}.
Berlin.
\url{https://dip.bundestag.de/\%C3\%BCber-dip/hilfe/api\#content}.

\bibitem[\citeproctext]{ref-schwarz_this_2026}
Schwarz, St. 2026. {``This Papers Corpus Build \& Evaluation Scripts.''}
\emph{GitHub/Esteeschwarz}. Berlin.
\url{https://github.com/esteeschwarz/SPUND-LX/blob/main/germanic/HA/}.

\bibitem[\citeproctext]{ref-wikipedia_google_2026}
Wikipedia, and Google. 2026. {``Google {Gemini}.''} \emph{Wikipedia}.
\url{https://de.wikipedia.org/w/index.php?title=Google_Gemini&oldid=263426206}.

\bibitem[\citeproctext]{ref-yakura_empirical_2025}
Yakura, Hiromu, Ezequiel Lopez-Lopez, Levin Brinkmann, Ignacio Serna,
Prateek Gupta, Ivan Soraperra, and Iyad Rahwan. 2025. {``Empirical
Evidence of {Large} {Language} {Model}'s Influence on Human Spoken
Communication.''} arXiv.
\url{https://doi.org/10.48550/arXiv.2409.01754}.

\end{CSLReferences}




\end{document}
